\documentclass[12pt]{article}
\usepackage[utf8]{inputenc}
\usepackage[T1]{fontenc}
\usepackage[margin=1in]{geometry}
\usepackage{setspace}
\usepackage{xcolor}
\usepackage[colorlinks=true, urlcolor=blue!60!black, linkcolor=blue!60!black]{hyperref}
\setlength{\parindent}{0em}
\setlength{\parskip}{0.6em}
\renewcommand{\baselinestretch}{1.0}
\begin{document}
\begin{center}
{\LARGE \textit{City of Memphis v.\ Greene}, 451 U.S.\ 100 (1981)\par}
\vspace{0.5em}
{\large Case Brief --- Two-Page Edition\par}
\end{center}

\section{Memory Jogger}
Closing a street between white and Black neighborhoods to cut through traffic did not violate 42 U.S.C.\ \S~1982 or the Thirteenth Amendment (\href{https://supreme.justia.com/cases/federal/us/451/100/}{Justia}).

\section{Detailed Case Facts}
Memphis closed the north end of West Drive, which ran through all-white Hein Park from a predominantly Black area, at the residents' request, citing through traffic and children's safety. About 1,000 opponents protested that the closure shut the door on Black neighbors.

\section{Procedural History}
After trial, the district court found disparate impact on ``overwhelmingly black'' traffic but no discriminatory intent. The Sixth Circuit reversed, calling the closure a ``badge of slavery.''

\section{Judicial Votes}
Reversed, 6--3. Stevens wrote for Burger, Stewart, Powell, and Rehnquist; White concurred in the judgment; Marshall dissented, joined by Brennan and Blackmun.

\section{Holding}
A neutral street closure serving legitimate traffic interests, without discriminatory intent, property depreciation, severe restriction of access, or unequal benefits, neither impairs \S~1982 property rights nor is a badge of slavery. The Court left open whether \S~1982 requires intent.

\section{Analysis of Opinions}
\textbf{Majority:} Drivers merely had to turn a few blocks earlier, property values would not fall, and the city denied no comparable benefit to Black neighborhoods. Ordinary traffic rules inevitably have disparate effects; calling them slavery would trivialize the Amendment. \textbf{White:} the Court should have decided that \S~1982 requires purposeful discrimination. \textbf{Marshall:} the barrier at the racial dividing line was a ``monument to racial hostility,'' backed by evidence of intent and procedural irregularity.

\section{Examples: Future Applications}
\textbf{Same-Side:} one-way conversion and speed humps on a white district's street to curb cut-through traffic. \textbf{Opposite-Side:} granting cul-de-sacs to white subdivisions while denying identical Black petitions. \textbf{Fence-Sitter:} a solid sound wall between a white subdivision and a Black apartment complex, prompted partly by complaints about ``foot traffic and crime.''

\section{Critique}
Critics say the Court ducked the intent question, echoed \textit{Plessy}'s dismissal of stigma, and froze Thirteenth Amendment litigation; textualists praise its restraint.

\section{Key Quotations}
\begin{quote}
``We merely hold that the impact of the closing of West Drive on nonresidents of Hein Park is a routine burden of citizenship; it does not reflect a violation of the Thirteenth Amendment.''
\end{quote}
\noindent --- Justice Stevens, for the Court, 451 U.S.\ at 129.

\end{document}
